\begin{comment}
  Homework solution of Calculus 2 - Chapter 4 Integral Application
  Licensed under MIT license
\end{comment}

\documentclass[12pt]{article}
\usepackage[utf8]{inputenc}
\usepackage{amsmath}
\usepackage{geometry}

\geometry{a4paper, margin=1in}

\title{No. 9 - Aplikasi Integral}
\author{Ahmad Fakhrul Bawani}
\date{}

\begin{document}

\maketitle

\section*{Solusi Volume Benda Putar (Metode Kulit Tabung)}

Diketahui daerah yang dibatasi oleh kurva $x = 3y^2$, sumbu-$x$ ($y=0$), dan garis $y = \sqrt{3}$ diputar mengelilingi sumbu-$x$. Karena pias diambil sejajar dengan sumbu putar, kita akan menggunakan metode kulit tabung (\textit{shell method}).

\subsection*{1. Analisis Komponen Kulit Tabung}
Saat diputar mengelilingi sumbu-$x$ dengan variabel integrasi $y$:
\begin{itemize}
  \item Jari-jari tabung ($r$) adalah jarak dari sumbu putar ke pias: $r = y$.
  \item Tinggi tabung ($h$) adalah panjang pias secara horizontal: $h = x = 3y^2$.
\end{itemize}

\subsection*{2. Rumus Volume Metode Kulit Tabung}
Volume ($V$) dihitung dengan mengintegralkan keliling dikali jari-jari dan tinggi terhadap $y$ dari $0$ hingga $\sqrt{3}$:
\[ V = \int_{c}^{d} 2\pi \cdot (\text{jari-jari}) \cdot (\text{tinggi}) \, dy \]
Substitusi nilai yang diketahui:
\[ V = \int_{0}^{\sqrt{3}} 2\pi \cdot y \cdot (3y^2) \, dy \]
\[ V = 2\pi \int_{0}^{\sqrt{3}} 3y^3 \, dy = 6\pi \int_{0}^{\sqrt{3}} y^3 \, dy \]

\subsection*{3. Menghitung Integral}
\begin{align*}
V &= 6\pi \left[ \frac{1}{4} y^4 \right]_{0}^{\sqrt{3}} \\
V &= \frac{6\pi}{4} \left[ (\sqrt{3})^4 - 0^4 \right] \\
V &= \frac{3\pi}{2} [ 9 ] \\
V &= \frac{27\pi}{2}
\end{align*}

\subsection*{Kesimpulan}
Jadi, volume benda putar yang dihasilkan menggunakan metode kulit tabung adalah $\mathbf{\frac{27\pi}{2}}$ atau \textbf{13,5$\pi$ satuan kubik}.

\end{document}